\part{Sheaf Theoretic Interpretations}
\label{part:XII}

Part XII interprets local knowledge and plans over overlapping contexts. Its stated payoff is the distinction between a presheaf of plans and a sheaf of feasible allocations, together with one genuine cohomological obstruction.

\begin{housing}Landlord, city, and contractor each hold local records, and gluing succeeds or fails depending on whether the shared crew is represented.\end{housing}

\chapter{Local Economic Knowledge and Overlapping Contexts}\label{ch:56}
Introduces a base of contexts such as jurisdictions, organizations, regions, or information domains. Restriction maps describe how a claim or plan appears within a narrower context. The choice of base is economically motivated and explicit.
\chspec{\Pref{P21}.}
\scaffold

\chapter{Presheaves of Plans, Accounts, and Commitments}\label{ch:57}
Local plans and records form presheaves when restriction is well defined. Asks whether accounts agree on overlaps, whether obligations receive consistent treatment across departments, and whether local descriptions omit shared dependencies.
\chspec{\Pref{P21}.}
\scaffold

\chapter{Gluing and Global Feasibility}\label{ch:58}
Interprets the sheaf condition as a compatibility-and-reconstruction property. Distinguishes agreement of descriptions from feasibility of the described activity: matching local plans can still compete for a globally scarce resource unless that resource is represented in the model. Examples show when gluing succeeds and when the proposed assignment is only a presheaf.
\chspec{\Pref{P09}, \Pref{P21}.}
\draftnote{Repair exercises: a local-repair condition need not produce a global section; a lower-energy imperfect candidate can beat an exact match under a penalized reconstruction; and a circle example must not state contradictory overlap values.}
\scaffold

\chapter{Obstructions, Inconsistency, and Institutional Boundaries}\label{ch:59}
Investigates failures of global coordination through carefully chosen mathematical models. Cohomological methods are introduced only where an appropriate coefficient structure and a justified obstruction interpretation exist. Distinguishes a genuine mathematical obstruction from a loose analogy about disagreement, beginning with the frictionless exchange graph.
\chspec{\Pref{P22}.}
\scaffold

\chapter{Descent, Federation, and the Organization of Collective Futures}\label{ch:60}
Examines how independently governed institutions can coordinate compatible records, rights, and obligations without one institution absorbing the others. Brings the three interpretations together: categorical composition for processes, type theoretic judgments for admissibility, and sheaf theoretic compatibility for local-to-global coordination.
\chspec{\Pref{P18}, \Pref{P20}, \Pref{P21}, \Pref{P22}.}
\scaffold
